\section{System Architecture on AWS}
\label{sec:arch}

\subsection{End-to-End Pipeline}
The pipeline keeps the three-phase shape of the previous semester (Figure~\ref{fig:arch}):

\begin{enumerate}
  \item \textbf{Phase 1 -- Extract.} English \code{.resx} files are converted to XLIFF, one file per target language,
        and uploaded to Amazon S3.
  \item \textbf{Phase 2 -- Translate.} Each XLIFF file is checked against the Translation Memory in Amazon RDS (SQL
        Server). Strings not found are translated by NLLB-200 (distilled, 600M parameters)~\cite{nllb} running on the
        CTranslate2 engine~\cite{ct2} with INT8 quantisation inside a Docker container, and the new translations are saved
        back to the memory.
  \item \textbf{Phase 3 -- Merge.} The translated XLIFF files are merged back into localized \code{.resx} files and
        published to S3.
\end{enumerate}

Every translated file emits one structured \code{RESEARCH\_DATA} telemetry record to Amazon CloudWatch Logs, with the
same schema as the previous semester plus the new fields listed in Appendix~\ref{app:telemetry}.

\begin{figure}[H]
\centering
\begin{tikzpicture}[
  box/.style={draw=cnew!70!black, fill=boxblue, rounded corners=3pt, align=center, font=\small, minimum height=1.25cm, text width=2.9cm},
  src/.style={draw=cgrey, fill=boxgrey, rounded corners=3pt, align=center, font=\small, minimum height=1.25cm, text width=2.3cm},
  store/.style={draw=cgrey, fill=boxgrey, rounded corners=3pt, align=center, font=\small, minimum height=1.25cm, text width=3.4cm},
  arr/.style={-{Stealth[length=2.4mm]}, thick, draw=black!70},
  lbl/.style={font=\scriptsize, fill=white, inner sep=1.5pt}
]
\node[box] (p1) {\textbf{Phase 1}\\Extract to XLIFF};
\node[box, right=0.9cm of p1] (p2) {\textbf{Phase 2}\\Translate with\\NLLB-200 on CT2};
\node[box, right=0.9cm of p2] (p3) {\textbf{Phase 3}\\Merge into \code{.resx}};
\node[src, left=1.3cm of p1] (commit) {Developer\\commit};
\node[store, below=1.6cm of p1] (s3) {\textbf{Amazon S3}\\XLIFF and outputs,\\model cache};
\node[store, below=1.6cm of p2] (rds) {\textbf{Amazon RDS}\\Translation Memory\\(SQL Server)};
\node[store, below=1.6cm of p3] (cw) {\textbf{CloudWatch Logs}\\logs and telemetry};
\node[store, above=1.3cm of p2] (sm) {\textbf{Secrets Manager}\\database credentials};
\begin{scope}[on background layer]
\node[draw=cnew!70!black, dashed, rounded corners=5pt, fit=(p1)(p2)(p3), inner sep=8pt] (runner) {};
\end{scope}
\node[font=\scriptsize\itshape, anchor=south west] at (runner.north west) {EC2 runner (m5.large, 2 vCPU, 8 GB)};
\draw[arr] (commit) -- node[lbl, above=3pt]{CI run} (p1.west -| runner.west);
\draw[arr] (p1) -- (p2);
\draw[arr] (p2) -- (p3);
\draw[arr, <->] (p1) -- node[lbl]{files} (s3);
\draw[arr, <->] (p2) -- node[lbl]{lookup / save} (rds);
\draw[arr] (p3) -- node[lbl]{telemetry} (cw);
\draw[arr] (sm) -- (sm |- runner.north);
\end{tikzpicture}
\caption[The localization pipeline on AWS]{The localization pipeline on AWS. All resources are created by one
CloudFormation stack in \code{ap-south-1}.}
\label{fig:arch}
\end{figure}

\subsection{Azure to AWS Service Mapping}
The move to AWS changed only the services around the translation code (Table~\ref{tab:mapping}). The translation logic
itself was kept identical to the Azure version, and automated parity tests compare the two code bases so that any
divergence is a deliberate, documented change.

\begin{table}[H]
\centering
\caption[Azure to AWS service mapping]{Azure services of the previous semester and their AWS replacements.}
\label{tab:mapping}
\begin{tabularx}{\linewidth}{@{}l l L@{}}
\toprule
\textbf{Role} & \textbf{Azure (previous semester)} & \textbf{AWS (this semester)} \\
\midrule
File storage & Blob Storage & Amazon S3 (\code{runs/<run-id>/raw}, \code{translated}, \code{output}) \\
Translation Memory & Azure SQL Database & Amazon RDS for SQL Server Express \\
Telemetry & Application Insights & Amazon CloudWatch Logs (same \code{RESEARCH\_DATA} payload) \\
Credentials & pipeline variables & AWS Secrets Manager (RDS-managed password) \\
CI/CD orchestration & Azure DevOps Pipelines & EC2 runner driven by SSM Run Command$^{\dagger}$ \\
Infrastructure & manual set-up & AWS CloudFormation (one stack) \\
\bottomrule
\end{tabularx}

\smallskip
{\footnotesize $^{\dagger}$An AWS CodePipeline/CodeBuild definition is included, but CodeBuild is not enabled for this
account, so runs use an EC2 runner with the same three phases.}
\end{table}

\subsection{Baseline on AWS}
The first complete AWS run (\code{ec2-20261015t233954z}) translated the previous semester's 50-string workload into seven
languages with an empty memory. All seven files completed with no failures (Table~\ref{tab:awsbaseline}). This run is the
starting point for the profiling in Chapter~\ref{sec:profiling}.

\begin{table}[H]
\centering
\caption[First AWS run]{First AWS run: 50 strings per language, empty Translation Memory, beam 5 (m5.large, 2 cores).}
\label{tab:awsbaseline}
\begin{tabular}{@{}l S[table-format=2.2] S[table-format=4.0] S[table-format=2.2]@{}}
\toprule
\textbf{Language} & {\textbf{File time (s)}} & {\textbf{Strings/hour}} & {\textbf{Model time (s)}} \\
\midrule
de-DE & 22.30 & 8070 & 22.15 \\
es-ES & 21.19 & 8492 & 21.04 \\
fr-FR & 22.73 & 7920 & 22.59 \\
hi-IN & 21.71 & 8291 & 21.57 \\
ja-JP & 20.41 & 8820 & 20.27 \\
ru-RU & 22.10 & 8143 & 21.97 \\
zh-CN & 19.47 & 9247 & 19.32 \\
\midrule
\textbf{Average} & \bfseries 21.42 & \bfseries 8426 & \bfseries 21.27 \\
\bottomrule
\end{tabular}
\end{table}
